\documentclass[a4paper,11pt, titlepage]{article}

%% Language and font encodings
\usepackage[english]{babel}
\usepackage[utf8]{inputenc}
\usepackage[T1]{fontenc}

%% Sets page size and margins
\usepackage[a4paper,top=1in,bottom=1in,left=1in,right=1in,marginparwidth=1.75cm]{geometry}

%% Useful packages
\usepackage{afterpage}
\usepackage{amsmath}
\usepackage{amsthm}
\usepackage{amssymb}
\usepackage{csquotes}
\usepackage{enumitem}
\usepackage{graphicx}
\usepackage{subcaption}
\usepackage{lipsum}
\usepackage{booktabs}
\newtheorem{proposition}{Proposition}
\theoremstyle{remark}
\newtheorem{remark}{Remark}
\usepackage{algorithm}
\usepackage{algpseudocode}
\usepackage[ruled,vlined]{algorithm2e}
\DontPrintSemicolon
\SetAlgoNoLine


% Listings (for displaying code):
\usepackage{listings}
\lstset{
    frame = single, 
    framexleftmargin=15pt
}

% Center figure captions:
\usepackage{caption}
\captionsetup[figure]{labelfont={bf},name={Figure},labelsep=quad}
\captionsetup[table]{labelfont={bf},name={Table},labelsep=quad}


% ----------- Algorithm2e setup
\usepackage[ruled,vlined]{algorithm2e}
\makeatletter
\renewcommand{\SetKwInOut}[2]{%
  \sbox\algocf@inoutbox{\KwSty{#2}\algocf@typo:}%
  \expandafter\ifx\csname InOutSizeDefined\endcsname\relax% if first time used
    \newcommand\InOutSizeDefined{}\setlength{\inoutsize}{\wd\algocf@inoutbox}%
    \sbox\algocf@inoutbox{\parbox[t]{\inoutsize}{\KwSty{#2}\algocf@typo:\hfill}~}\setlength{\inoutindent}{\wd\algocf@inoutbox}%
  \else% else keep the larger dimension
    \ifdim\wd\algocf@inoutbox>\inoutsize%
    \setlength{\inoutsize}{\wd\algocf@inoutbox}%
    \sbox\algocf@inoutbox{\parbox[t]{\inoutsize}{\KwSty{#2}\algocf@typo:\hfill}~}\setlength{\inoutindent}{\wd\algocf@inoutbox}%
    \fi%
  \fi% the dimension of the box is now defined.
  \algocf@newcommand{#1}[1]{%
    \ifthenelse{\boolean{algocf@inoutnumbered}}{\relax}{\everypar={\relax}}%
%     {\let\\\algocf@newinout\hangindent=\wd\algocf@inoutbox\hangafter=1\parbox[t]{\inoutsize}{\KwSty{#2}\algocf@typo\hfill:}~##1\par}%
    {\let\\\algocf@newinout\hangindent=\inoutindent\hangafter=1\parbox[t]{\inoutsize}{\KwSty{#2}\algocf@typo:\hfill}~##1\par}%
    \algocf@linesnumbered% reset the numbering of the lines
  }}%
\makeatother
% --------- end algorithm2e setup

% \bm allows typing bold math:
\usepackage{bm}
\usepackage[normalem]{ulem}

\usepackage[colorinlistoftodos]{todonotes}
\usepackage[colorlinks=true, allcolors=blue]{hyperref}

\renewcommand*{\rmdefault}{bch}
\renewcommand*{\ttdefault}{lmtt}
\newcommand{\citationneeded}{\textcolor{red}{[citation-needed]}}

\DeclareMathOperator*{\argmin}{\arg\!\min}
\DeclareMathOperator*{\argmax}{\arg\!\max}

% Add bigger skip between paragraphs, makes reading easier:
\setlength{\parskip}{0.5em}

% Bibliography
\usepackage[numbers, comma, square, sort&compress]{natbib}
\bibliographystyle{abbrvunsrtnat.bst}


% Theorem, Lemma, Corollary and Proofs
\theoremstyle{definition}
\newtheorem{definition}{Definition}[section]
\theoremstyle{plain}
\newtheorem{theorem}{Theorem}[section]
\newtheorem{corollary}{Corollary}[theorem]
\newtheorem{lemma}[theorem]{Lemma}
\theoremstyle{remark}

%----------------------------------------------------------------------------------------
% END OF DOCUMENT CONFIGURATION
%----------------------------------------------------------------------------------------

%----------------------------------------------------------------------------------------
% IMPORTANT INFORMATION TO MODIFY FOR THE TITLE PAGE
%----------------------------------------------------------------------------------------
\newcommand{\reporttitle}{Title} % Title of your research project
\newcommand{\reportauthorA}{Student name 1 (CID: -------------)} % First Name, Last Name and CID of student 1

\begin{document}

%%%%%%%%%%%%%%%%%%%%%%%%%%%%%%%%%%%%%
% Title page
%%%%%%%%%%%%%%%%%%%%%%%%%%%%%%%%%%%%%
\begin{titlepage}
\newcommand{\HRule}{\rule{\linewidth}{0.5mm}} % Defines a new command for the horizontal lines, change thickness here
%----------------------------------------------------------------------------------------
%	LOGO SECTION
%----------------------------------------------------------------------------------------
\includegraphics[width=8cm]{Imperial_logo.png}\\[1cm] % Include a department/university logo - this will require the graphicx package
%----------------------------------------------------------------------------------------
\center % Center everything on the page
%----------------------------------------------------------------------------------------
%	HEADING SECTIONS
%----------------------------------------------------------------------------------------
% For M3R or M4R reports, comment out as appropriate
\textsc{\LARGE Imperial College London}\\[0.5cm] 
\textsc{\Large Department of Mathematics}\\[1.5cm] 
\textsc{\Large Second-year Group Research Project}\\[0.5cm]
%----------------------------------------------------------------------------------------
%	TITLE SECTION
%----------------------------------------------------------------------------------------
\makeatletter
\HRule \\[0.6cm]
{ \huge \bfseries \reporttitle}\\[0.6cm] % Title of your document
\HRule \\[1.5cm]
%----------------------------------------------------------------------------------------
%	AUTHOR SECTION
%----------------------------------------------------------------------------------------
\begin{minipage}{0.4\textwidth}
\begin{flushleft} \large
\emph{Author:}\\
\reportauthorA \\
\reportauthorB \\
\reportauthorC \\
\reportauthorD \\
\reportauthorE
\end{flushleft}
\end{minipage}
~
\begin{minipage}{0.4\textwidth}
\begin{flushright} \large
\emph{Supervisor(s):} \\
\supervisor % Supervisor's name
\end{flushright}
\end{minipage}\\[2cm]
\makeatother
%----------------------------------------------------------------------------------------
%	FOOTER & DATE SECTION
%----------------------------------------------------------------------------------------
\vfill
\makeatletter
{\large \today}\\[2cm] % Date, change the \today to a set date if you want to be precise
\makeatother
\end{titlepage}
%%%%%%%%%%%%%%%%%%%%%%%%%%%%%%%%%%%%%


%%%%%%%%%%%%%%%%%%%%%%%%%%%%%%%%%%%%%
% Table of contents
\tableofcontents
\clearpage
%%%%%%%%%%%%%%%%%%%%%%%%%%%%%%%%%%%%%

%%%%%%%%%%%%%%%%%%%%%%%%%%%%%%%%%%%%%
% Main sections
%%%%%%%%%%%%%%%%%%%%%%%%%%%%%%%%%%%%%

%%%%%%%%%%%%%%%%%%%%%%%%%%%%%%%%%%%%%
% Main sections
%%%%%%%%%%%%%%%%%%%%%%%%%%%%%%%%%%%%%

\section{Introduction}
\label{sec:introduction}

A discrete dynamical system is generated by iterating a map. Although the rule may be simple, the long-term behaviour of its orbits can be highly complicated. In this report, the main object is the Hénon map with the classical parameter choice and the Hénon attractor.

The main theme of the report is hyperbolic behaviour. The Smale horseshoe provides the clean reference model: stretching and folding lead to symbolic dynamics, sensitive dependence on initial conditions, and a uniformly hyperbolic invariant set. The Hénon attractor has a similar geometric mechanism, but the hyperbolic structure is more complicated. 

The numerical part of this report studies typical hyperbolic behaviour along long orbits. The relevant theoretical background comes from invariant measures, Birkhoff's Ergodic Theorem, and Oseledets' Multiplicative Ergodic Theorem, but the computations in this report remain finite-time diagnostics rather than asymptotic proofs. The computations support a hyperbolic-type pattern.

\section{Theoretical background}
\label{sec:theoretical-model}

\subsection{Discrete dynamical systems and orbits}
\label{subsec:discrete-dynamical-systems}

\begin{definition}[Discrete dynamical system]
Let \(X\) be a metric space and let \(f:X\to X\) be a map. The discrete dynamical system generated by \(f\) is the iteration
\[
x_{n+1}=f(x_n), \qquad x_n=f^n(x_0).
\]
\end{definition}

\begin{definition}[Forward orbit]
For \(x_0\in X\), the forward orbit of \(x_0\) is
\[
\mathcal O^+(x_0)
=
\{x_0,f(x_0),f^2(x_0),\ldots\}.
\]
\end{definition}

\subsection{Invariant sets and attractors}
\label{subsec:invariant-sets}

We now recall the basic terminology for invariant and attracting sets, following Robinson's notation. Let \(f:M\to M\) be a diffeomorphism.

\begin{definition}[Invariant set]
A set \(\Lambda\subset M\) is called invariant under \(f\) if
\(f(\Lambda)=\Lambda .\)
It is called positively invariant if
\(f(\Lambda)\subset \Lambda .\)
\end{definition}

Thus an invariant set is a dynamically closed set on which the long-term dynamics can be studied.

\begin{definition}[Trapping region]
A compact region \(N\subset M\) is called a trapping region for \(f\) if
\[
f(N)\subset \operatorname{int}(N).
\]
\end{definition}

The forward images of a trapping region form a nested sequence
\[
N\supset f(N)\supset f^2(N)\supset \cdots .
\]

\begin{definition}[Attracting set]
If \(N\) is a trapping region, the associated attracting set is
\[
\Lambda=\bigcap_{k\geq 0} f^k(N).
\]
\end{definition}

The attracting set is the limiting set left after repeatedly applying the map to the trapping region. It is invariant, since
\[
f(\Lambda)
=
f\left(\bigcap_{k\geq 0}f^k(N)\right)
=
\bigcap_{k\geq 1}f^k(N)
=
\Lambda .
\]

\begin{definition}[Basin of attraction]
The basin of attraction of an attracting set \(\Lambda\) is the set of points whose forward orbits approach (\Lambda):
\[
W^s(\Lambda)
=
{x\in M:\operatorname{dist}(f^n(x),\Lambda)\to 0
\text{ as } n\to\infty}.
\]
\end{definition}

So \(\Lambda\) is the set being approached, while \(W^s(\Lambda)\) is the set of initial conditions attracted to it.

\begin{definition}[\(\varepsilon\)-chain and chain transitivity]
Let \(\varepsilon>0\). An \(\varepsilon\)-chain from \(x\) to \(y\) in \(\Lambda\) is a finite sequence
\[
x=x_0,x_1,\ldots,x_m=y
\]
such that
\[
d(f(x_i),x_{i+1})<\varepsilon
\qquad\text{for } i=0,\ldots,m-1.
\]
The map \(f|_{\Lambda}\) is called chain transitive if, for any \(x,y\in\Lambda\) and any \(\varepsilon>0\), there exists an \(\varepsilon\)-chain from \(x\) to \(y\) contained in \(\Lambda\).
\end{definition}

\begin{definition}[Attractor]
An attracting set \(\Lambda\) is called an attractor if the restricted dynamics \(f|_{\Lambda}\) is chain transitive.
\end{definition}

For the Hénon map \(F_{AB}\), the numerically observed strange attractor is treated as a finite approximation of an attracting invariant set. In practice, we iterate a typical initial condition for a long time, discard the transient part, and use the remaining orbit points as a sample of the observed attracting set.

\subsection{Tangent space and linearisation}
\label{subsec:Tangent space and linearisation}

\begin{definition}[Tangent vector and tangent space in Euclidean space]
Let \(p\in \mathbb R^n\). 

\begin{enumerate}[label=(\arabic*)]
    \item A tangent vector at \(p\) is a pair \(v_p = (p, v)\), $v\in \mathbb R^n$. The point \(p\) is called the base point of the tangent vector, and \(v\) is called its direction vector. 

    \item The tangent space at \(p\), denoted by \(T_p\mathbb R^n\), is the set of all
    tangent vectors based at \(p\):
    \[
    T_p\mathbb R^n
    =
    \{(p,v):v\in\mathbb R^n\}.
    \]
    It is a vector space under the operations
    \[
    (p,v)+(p,w)=(p,v+w),
    \]
    and
    \[
    \alpha(p,v)=(p,\alpha v),
    \qquad \alpha\in\mathbb R.
    \]
    Hence \(T_p\mathbb R^n\) is isomorphic to \(\mathbb R^n\).
    
    \item The tangent bundle/space of \(\mathbb R^n\) is the disjoint union of the tangent
    spaces at all points:
    \[
    T\mathbb R^n
    =
    \bigcup_{p\in\mathbb R^n}T_p\mathbb R^n.
    \]
    Equivalently,
    \[
    T\mathbb R^n
    =
    \{(p,v):p\in\mathbb R^n,\ v\in\mathbb R^n\}
    =
    \mathbb R^n\times\mathbb R^n.
    \]
\end{enumerate}

\end{definition}

\label{subsec:linearisation}

\begin{definition}[Derivative and linearisation in Euclidean space]
Let \(f:\mathbb R^n\to\mathbb R^m\) be a \(C^1\) map. The derivative of \(f\)
at \(p\in\mathbb R^n\) is the linear map
\[
Df_p:T_p\mathbb R^n\to T_{f(p)}\mathbb R^m
\]
defined by
\[
Df_p(p,v)=(f(p),Df(p)v),
\]
where \(Df(p)\) is the Jacobian matrix of \(f\) at \(p\).

The derivative \(Df_p\) is also called the linearisation of \(f\) at \(p\),
because it gives the first-order approximation
\[
f(p+v)=f(p)+Df(p)v+o(\|v\|)
\qquad \text{as } \|v\|\to 0.
\]
\end{definition}

\begin{definition}[Derivative cocycle]
Let \(f:\mathbb R^n\to\mathbb R^n\) be a \(C^1\) map. For \(k\geq 1\), define
\[
Df_p^k := D(f^k)_p : T_p\mathbb R^n \to T_{f^k(p)}\mathbb R^n .
\]
By the chain rule,
\[
Df_p^k
=
Df_{f^{k-1}(p)}
Df_{f^{k-2}(p)}
\cdots
Df_{f(p)}
Df_p .
\]
Equivalently, if \(p_j=f^j(p)\), then
\[
Df_p^k
=
Df_{p_{k-1}}Df_{p_{k-2}}\cdots Df_{p_1}Df_{p_0}.
\]

The family
\[
(p,k)\mapsto Df_p^k
\]
is called the derivative cocycle over \(f\). It satisfies the cocycle identity
\[
Df_p^{k+\ell}
=
Df_{f^\ell(p)}^k Df_p^\ell
\qquad
\text{for all } k,\ell\geq 0.
\]
\end{definition}

Let \(p_0\in\mathbb R^n\) and consider its orbit
\[
p_{j+1}=f(p_j),
\qquad j\geq 0.
\]
The derivative cocycle describes the evolution of infinitesimal perturbations
along an orbit. Let
\[
p_0=p,
\qquad
p_{j+1}=f(p_j)=f^{j+1}(p).
\]
If \(v_0\in T_{p_0}\mathbb R^n\cong\mathbb R^n\) is a small perturbation of
the initial point \(p_0\), then after one step the perturbation is approximated by
\[
v_1=Df(p_0)v_0.
\]
After two steps,
\[
v_2=Df(p_1)v_1
=
Df(p_1)Df(p_0)v_0.
\]
Inductively, after \(k\) steps,
\[
v_k
=
Df(p_{k-1})\cdots Df(p_1)Df(p_0)v_0
=
A(k,p_0)v_0.
\]
Thus the nonlinear orbit \(p_{j+1}=f(p_j)\) induces the linearised perturbation dynamics by
\(
v_{j+1}=Df(p_j)v_j.
\)




\subsection{Hyperbolicity}
\label{subsec:Hyperbolicity}
\begin{definition}[Hyperbolic structure and hyperbolic invariant set]
Let \(f:\mathbb R^n\to\mathbb R^n\) be a \(C^1\) diffeomorphism, and let
\(\Lambda\subset \mathbb R^n\) be a compact invariant set. We say that
\(\Lambda\) has a hyperbolic structure for \(f\), or \(\Lambda\) is a hyperbolic invariant set if for every \(p\in\Lambda\),
the tangent space \(T_p\mathbb R^n\) admits a direct-sum decomposition
\[
T_p\mathbb R^n
=
E_p^s\oplus E_p^u,
\]
such that:

\begin{enumerate}[label=(\arabic*)]
    \item The splitting is invariant under the derivative:
    \[
    Df_p(E_p^s)=E_{f(p)}^s,
    \qquad
    Df_p(E_p^u)=E_{f(p)}^u.
    \]

    \item There exist constants \(C\geq 1\) and \(0<\lambda<1\), independent of
    \(p\in\Lambda\), such that for all \(k\geq 0\),
    \[
    \|Df_p^k v^s\|
    \leq
    C\lambda^k\|v^s\|,
    \qquad
    v^s\in E_p^s,
    \]
    and
    \[
    \|Df_p^{-k} v^u\|
    \leq
    C\lambda^k\|v^u\|,
    \qquad
    v^u\in E_p^u.
    \]
\end{enumerate}

The subspaces \(E_p^s\) and \(E_p^u\) are called the stable and unstable
subspaces at \(p\), respectively.
\end{definition}

\begin{remark}[Uniformity in the definition]
In the above definition, the constants \(C\geq 1\) and \(0<\lambda<1\) are
independent of the base point \(p\in\Lambda\). Therefore a hyperbolic invariant
set in this sense is already uniformly hyperbolic.

This is different from nonuniform hyperbolicity, where one typically assumes
that the Lyapunov exponents are nonzero for almost every point with respect to
an invariant measure. In that setting, the stable and unstable splitting may exist
only almost everywhere, may be merely measurable rather than continuous, and
the contraction and expansion rates need not be controlled by constants uniform
over the whole set[].
\end{remark}

\begin{definition}[Stable and unstable directions in dimension two]
Suppose \(f:\mathbb R^2\to\mathbb R^2\) is a \(C^1\) diffeomorphism and
\(\Lambda\subset\mathbb R^2\) is a hyperbolic invariant set. If for each \(p\in\Lambda\),
\[
\dim E_p^s=\dim E_p^u=1,
\]
then \(E_p^s\) and \(E_p^u\) are called the stable and unstable directions at
\(p\), respectively.

In this case, \(E_p^s\) and \(E_p^u\) are one-dimensional subspaces of
\(
T_p\mathbb R^2\cong \mathbb R^2.
\)
\end{definition}

\begin{remark}[Meaning of the stable and unstable subspaces]
The stable subspace \(E_p^s\) consists of infinitesimal directions that are
contracted exponentially under forward iteration. The unstable subspace
\(E_p^u\) consists of infinitesimal directions that are contracted exponentially
under backward iteration, equivalently expanded exponentially under forward
iteration.

More precisely, if \(v^s\in E_p^s\), then
\[
\|Df_p^k v^s\|
\leq
C\lambda^k\|v^s\|
\qquad
\text{as } k\to\infty.
\]
If \(v^u\in E_p^u\), then
\[
\|Df_p^{-k} v^u\|
\leq
C\lambda^k\|v^u\|
\qquad
\text{as } k\to\infty.
\]
The use of backward iterates in the definition of \(E_p^u\) is important:
forward iteration expands unstable vectors, so backward iteration is the
contracting direction used to characterize them.
\end{remark}

\begin{definition}[Angle between stable and unstable directions] Suppose \(f:\mathbb R^2\to\mathbb R^2\) is a \(C^1\) diffeomorphism and \(\Lambda\subset\mathbb R^2\) is a hyperbolic invariant set with one-dimensional stable and unstable directions. Let \(e_p^s\) and \(e_p^u\) be unit vectors spanning \(E_p^s\) and \(E_p^u\), respectively. The angle between the stable and unstable directions at \(p\) is defined by \[ \theta(p)=\angle(E_p^s,E_p^u)= \arccos\left( \left|\langle e_p^s,e_p^u\rangle\right| \right) \in [0,\pi/2]. \] \end{definition} 

\begin{remark}[Meaning of the unit vectors \(e_p^s\) and \(e_p^u\)]
The vectors \(e_p^s\) and \(e_p^u\) are unit vectors chosen to span the
one-dimensional stable and unstable subspaces:
\[
E_p^s=\operatorname{span}(e_p^s),
\qquad
E_p^u=\operatorname{span}(e_p^u).
\]
They are not, in general, eigenvectors of the single matrix \(Df(p)\). The
reason is that
\[
Df_p:T_p\mathbb R^2\to T_{f(p)}\mathbb R^2
\]
maps tangent vectors at \(p\) to tangent vectors at \(f(p)\), rather than mapping
\(T_p\mathbb R^2\) to itself. The invariant splitting is instead characterised by
\[
Df_p(E_p^s)=E_{f(p)}^s,
\qquad
Df_p(E_p^u)=E_{f(p)}^u.
\]
Thus the stable and unstable directions are directions transported along the
orbit by the derivative cocycle.

At a fixed point \(p=f(p)\), these directions reduce to the stable and unstable
eigenspaces of \(Df(p)\). For a general point on a hyperbolic invariant set, they
are not pointwise eigenvectors of \(Df(p)\), but invariant directions along the
orbit.
\end{remark}

\begin{proposition}[Basic properties of hyperbolic sets]
Let \(\Lambda\) be a hyperbolic invariant set for a \(C^1\) diffeomorphism
\(f:\mathbb R^n\to\mathbb R^n\). Then the following properties hold:

\begin{enumerate}[label=(\arabic*)]
    \item The stable and unstable subspaces vary continuously with the base point:
    \[
    p\mapsto E_p^s,
    \qquad
    p\mapsto E_p^u.
    \]

    \item The stable and unstable subspaces are uniformly transverse:
    \[
    \inf_{p\in\Lambda}
    \angle(E_p^s,E_p^u)
    >
    0.
    \]
\end{enumerate}
\end{proposition}

The first property shows if two points \(p,q\in\Lambda\) are close, then the corresponding subspaces \(E_p^s,E_q^s\) and \(E_p^u,E_q^u\) are also close. Thus the stable and unstable directions do not jump discontinuously as the base point moves along \(\Lambda\). 

The second property says that the stable and unstable subspaces remain uniformly separated. In dimension two, this means that the stable and unstable directions are never tangent, and their angle is bounded away from zero. This is stronger than merely requiring \(E_p^s\neq E_p^u\) at each point. It rules out the possibility that the two directions become arbitrarily close somewhere on the invariant set. 

These two properties are important for the numerical diagnostics used later. Continuity means that stable and unstable directions can be meaningfully tracked along an orbit. Uniform transversality means that, for a uniformly hyperbolic set, the angle between the stable and unstable directions should not approach zero. Therefore, observing very small angles in finite-time computations can be interpreted as evidence of non-uniform behaviour or a possible near-tangency.


\subsection{The Smale horseshoe as a reference model}
\label{subsec:Smale-horseshoe}

The Smale horseshoe is a standard example of a uniformly hyperbolic invariant
set with complicated dynamics. It provides a clean reference case
where hyperbolicity and symbolic dynamics can be described explicitly.

\begin{definition}[Full two-shift]
Let
\[
\Sigma_2=\{0,1\}^{\mathbb Z}
\]
be the space of bi-infinite sequences on two symbols. The shift map
\[
\sigma:\Sigma_2\to\Sigma_2
\]
is defined by
\[
(\sigma s)_j=s_{j+1},
\qquad
s=(s_j)_{j\in\mathbb Z}\in\Sigma_2.
\]
\end{definition}

\begin{definition}[Horseshoe invariant set]
Let \(R\subset\mathbb R^2\) be the rectangle used in the horseshoe construction,
and let \(f:\mathbb R^2\to\mathbb R^2\) be the corresponding horseshoe
diffeomorphism. The invariant set of the horseshoe is
\[
\Lambda
=
\bigcap_{n\in\mathbb Z} f^n(R).
\]
Thus \(\Lambda\) consists of the points whose full forward and backward orbits
remain in \(R\).
\end{definition}

\begin{theorem}[Basic properties of the Smale horseshoe]
For the Smale horseshoe map \(f\), the invariant set
\[
\Lambda
=
\bigcap_{n\in\mathbb Z} f^n(R)
\]
has the following properties:

\begin{enumerate}[label=(\arabic*)]
    \item \(\Lambda\) is compact and invariant:
    \[
    f(\Lambda)=\Lambda.
    \]

    \item \(\Lambda\) is a Cantor-type set.

    \item \(\Lambda\) is uniformly hyperbolic. In particular, for each
    \(p\in\Lambda\),
    \[
    T_p\mathbb R^2
    =
    E_p^s\oplus E_p^u,
    \]
    where \(E_p^s\) is contracted under forward iteration and \(E_p^u\) is
    expanded under forward iteration.

    \item The dynamics of \(f\) on \(\Lambda\) is topologically conjugate to the
    full two-shift. That is, there exists a homeomorphism
    \[
    h:\Lambda\to\Sigma_2
    \]
    such that
    \[
    h\circ f=\sigma\circ h.
    \]
\end{enumerate}
\end{theorem}

The horseshoe provides the ideal uniformly hyperbolic reference case. We now
turn to the Hénon map, which has a similar geometric mechanism of stretching,
folding, and dissipation, but whose classical attractor is much less rigid.





\section{The Hénon map}
\label{sec:Hénon-map}

\subsection{Definition and basic properties}
\label{Definition and basic properties}





\begin{definition}[Hénon map]
For parameters \(A,B\in\mathbb R\), define
\[
F_{AB}(x,y)
=
(A-By-x^2,x).
\]
The classical parameters are $A=1.4, B=-0.3.$
\end{definition}

For the classical parameters \(A=1.4\) and \(B=-0.3\), numerical iteration
suggests that long orbits approach a thin attracting set of this type, usually
called the Hénon attractor. []


\begin{proposition}
\begin{enumerate}[label=(\arabic*)]
    \item The derivative of the Hénon map is
    \[
    DF_{AB}(x,y)
    =
    \begin{pmatrix}
    -2x & -B\\
    1 & 0
    \end{pmatrix}.
    \]
    
    Moreover,
    \[
    \det DF_{AB}(x,y)=B.
    \]
    \item If \(B\neq 0\), then \(F_{AB}\) is invertible and
    \[
    F^{-1}_{AB}(x,y)
    =
    \left(
    y,\frac{A-x-y^2}{B}
    \right).
    \]
\end{enumerate}

\end{proposition}

\begin{remark}[Equivalent with $H(a,b)$]
    
\end{remark}

\subsection{Geometric mechanism: stretching, folding, and dissipation}
\label{subsec:Geometric mechanism}

The Hénon map is useful because it gives a simple two-dimensional model of the
geometric mechanism behind many chaotic attractors. It combines three effects:
stretching, folding, and dissipation.

\begin{enumerate}
    \item The map is nonlinear in the \(x\)-coordinate. The term \(-x^2\) bends the image of vertical and horizontal sets and causes folding.
    \item The derivative
    \[
    DF_{AB}(x,y)
    =
    \begin{pmatrix}
    -2x & -B\\
    1 & 0
    \end{pmatrix}
    \]
    depends on the current position, so the local linear behaviour changes along
    the orbit. The derivative cocycle
    \(
    D(F_{AB}^k)(z_0)
    =
    DF_{AB}(z_{k-1})\cdots DF_{AB}(z_0)
    \)
    records the cumulative stretching, contraction, and rotation of infinitesimal
    directions along the orbit.

    \item The determinant is constant:
    \(
    \det DF_{AB}(x,y)=B.
    \)
    When $|B|<1$, the map is dissipative. For the classical parameter choice, the areas are contracted by the factor \(|B|=0.3\) at each iterate, and the negative
    sign indicates orientation reversal.
\end{enumerate}


The combination of these effects is the essential mechanism. Depending on the parameters, the maximum invariant set can have different geometric forms. For fixed $B$, sufficiently large $A$ produces a horseshoe-type structure, while some A, such as the classical parameter, produces a strange attractor. This makes the Hénon map a natural object for studying the relation between
Lyapunov exponents, stable--unstable directions, and fractal geometry.




\section{Numerical diagnostics}
\label{sec:numerical-diagnostics}

The numerical part of this report studies the linearised dynamics along long
orbits of the Hénon map. The basic objects are derivative products, their singular
values, finite-time Lyapunov exponents, and finite-time approximations of stable
and unstable directions.

\subsection{Theoretical basis}
\label{Theoretical basis}

\subsubsection{Derivative products along an orbit}

Let
\(F=F_{AB}:\mathbb R^2\to\mathbb R^2\)
be the Hénon map, and let \(z_{j+1}=F(z_j)\)
be an orbit. The derivative cocycle over this orbit is
\[
A_n(z_0)
=
D(F^n)(z_0)
=
DF(z_{n-1})DF(z_{n-2})\cdots DF(z_0).
\]
Thus a tangent perturbation \(v_0\in T_{z_0}\mathbb R^2\cong\mathbb R^2\)
evolves according to $v_n=A_n(z_0)v_0.$

The product \(A_n(z_0)\) records the cumulative stretching, contraction, and
rotation of infinitesimal tangent directions along the orbit segment $z_0,z_1,\dots,z_n.$

\subsubsection{Singular values and finite-time Lyapunov exponents}

Let \(A\in\mathbb R^{2\times 2}\). The singular values of \(A\) are \(\sigma_1(A)\geq \sigma_2(A)\geq 0,
\)
where
\[
\sigma_i(A)
=
\sqrt{\lambda_i(A^T A)},\qquad i=1,2.
\]
Here \(\lambda_i(A^T A)\) denotes the \(i\)-th eigenvalue of the symmetric
positive semidefinite matrix \(A^T A\).

For the derivative product
\(
A_n(z)=D(F^n)(z),
\)
the finite-time Lyapunov exponents over \(n\) iterates are defined by
\[
\lambda_i^{(n)}(z)
=
\frac{1}{n}\log \sigma_i(A_n(z))=
\frac{1}{2n}
\log \lambda_i\left(A_n(z)^T A_n(z)\right),
\qquad
i=1,2.
\]

The largest finite-time Lyapunov exponent is therefore
\[
\lambda_1^{(n)}(z)
=
\frac{1}{n}\log \|A_n(z)\|,
\]
where \(\|\cdot\|\) is the operator norm induced by the Euclidean norm.

\begin{remark}[Area contraction for the Hénon map]
Since
\(
\det DF_{AB}(x,y)=B,
\)
we have \(
\det A_n(z)=B^n.
\)
Hence
\[
\sigma_1(A_n(z))\sigma_2(A_n(z))
=
|\det A_n(z)|
=
|B|^n.
\]
Therefore,
\[
\lambda_1^{(n)}(z)+\lambda_2^{(n)}(z)
=
\log |B|.
\]
For the classical parameters, this gives\(
\lambda_1^{(n)}(z)+\lambda_2^{(n)}(z)
=
\log 0.3<0\), reflecting the dissipative nature of the Hénon map.

In the numerical computations, we compute \(\lambda_1^{(n)}(z)\) from the
largest singular value and then obtain
\(
\lambda_2^{(n)}(z)
=
\log|B|-\lambda_1^{(n)}(z).
\)
\end{remark}



\subsubsection{Oseledets' theorem}

The theoretical basis for using derivative products and singular values is the
Multiplicative Ergodic Theorem, also called Oseledets' theorem.

\begin{theorem}[Multiplicative Ergodic Theorem for a \(C^1\) diffeomorphism]
Let \(F:\mathbb R^2\to\mathbb R^2\) be a \(C^1\) diffeomorphism, let
\(\Lambda\subset\mathbb R^2\) be a compact invariant set, and let \(\mu\) be an
ergodic \(F\)-invariant probability measure supported on \(\Lambda\). For
\(n\geq 1\), define the derivative cocycle
\(
A_n(z)=D(F^n)(z).
\)
Then there exists a set \(\Lambda_0\subset\Lambda\) with
\[
\mu(\Lambda_0)=1,
\qquad
F(\Lambda_0)=\Lambda_0,
\]
such that for every \(z\in\Lambda_0\), the following properties hold.

\begin{enumerate}[label=(\arabic*)]
    \item There exist Lyapunov exponents
    \(
    \lambda_1>\lambda_2>\cdots>\lambda_\ell
    \)
    and a measurable splitting
    \[
    T_z\mathbb R^2
    =
    E_1(z)\oplus E_2(z)\oplus\cdots\oplus E_\ell(z),
    \]
    such that the splitting is invariant under the derivative:
    \[
    DF(z)E_i(z)=E_i(F(z)),
    \qquad
    i=1,\dots,\ell.
    \]
    
    \item For every nonzero vector \(v\in E_i(z)\),
    \[
    \lim_{n\to\infty}
    \frac{1}{n}\log \|A_n(z)v\|
    =
    \lambda_i.
    \]
    Equivalently, if
    \[
    \sigma_1(A_n(z))\geq \sigma_2(A_n(z))>0
    \]
    are the singular values of \(A_n(z)\), then the limits
    \[
    \lim_{n\to\infty}
    \frac{1}{n}\log \sigma_i(A_n(z))
    \]
    exist for \(i=1,2\). These limits are the Lyapunov exponents counted with
    multiplicity.
    
\end{enumerate}


\end{theorem}


The theorem shows that asymptotic Lyapunov exponents can be obtained from
singular values of the derivative cocycle:
\[
\lambda_i
=
\lim_{n\to\infty}
\frac{1}{n}\log \sigma_i(D(F^n)(z)).
\]
This is why the finite-time Lyapunov exponents used in the numerical part are
defined by
\[
\lambda_i^{(n)}(z)
=
\frac{1}{n}\log \sigma_i(D(F^n)(z)).
\]
They are finite-window approximations of the asymptotic limits given by
Oseledets' theorem.

\begin{remark}[Relation with stable and unstable directions]
Oseledets' theorem gives measurable subspaces \(E_i(z)\) on a full-measure set.
When there are two distinct Lyapunov exponents
\[
\lambda^u>\lambda^s,
\]
one obtains a splitting
\[
T_z\mathbb R^2
=
E^u(z)\oplus E^s(z),
\]
where vectors in \(E^u(z)\) grow at rate \(\lambda^u\), and vectors in
\(E^s(z)\) grow at rate \(\lambda^s\).

For a uniformly hyperbolic invariant set, this measurable Oseledets splitting
agrees almost everywhere with the continuous hyperbolic splitting
\[
T_z\mathbb R^2=E_z^s\oplus E_z^u.
\]
For a nonuniformly hyperbolic system, the Oseledets splitting may exist only
almost everywhere and may be merely measurable rather than continuous.
\end{remark}

\begin{remark}[Finite-time interpretation]
In numerical computations, the limit \(n\to\infty\) cannot be taken. Instead,
one studies
\[
\lambda_i^{(n)}(z)
=
\frac{1}{n}\log \sigma_i(A_n(z)).
\]
The dependence of \(\lambda_i^{(n)}(z)\) on both \(z\) and \(n\) measures how
finite-time stretching varies along the sampled orbit. Strong variation of these
finite-time quantities indicates non-uniformity of expansion and contraction
over the observed time windows.
\end{remark}

\subsubsection{Finite-time stable and unstable directions}

For a hyperbolic invariant set, the tangent space splits as $T_z\mathbb R^2=E_z^s\oplus E_z^u$, thus any tangent vector \(v\in T_z\mathbb R^2\) can be written uniquely as
\[
v=v^s+v^u,
\qquad
v^s\in E_z^s,\quad v^u\in E_z^u.
\]
Under forward iteration,
\[
DF_z^k v
=
DF_z^k v^s+DF_z^k v^u.
\]
The stable component \(DF_z^k v^s\) is contracted, while the unstable component
\(DF_z^k v^u\) is expanded. Therefore, unless \(v\in E_z^s\), the direction of
\(DF_z^k v\) aligns with the unstable direction as \(k\to\infty\). Similarly,
under backward iteration, a generic vector aligns with the stable direction.

Thus, numerically, \(E_z^u\) is approximated by forward iteration from the
past, and \(E_z^s\) is approximated by backward iteration from the future.

\begin{definition}[Finite-time unstable and stable direction] Let \(N\geq 1\).
\begin{enumerate}[label=(\arabic*)]
    \item Let \(z_{-N}\) be a point on the backward orbit of \(z_0\),
    so that \(F^N(z_{-N})=z_0.\)
    Choose any nonzero vector \(u_{-N}\in T_{z_{-N}}\mathbb R^2\) which is not in
    \(E_{z_{-N}}^s\), and define
    \[
    u_{j+1}
    =
    \frac{DF(z_j)u_j}{\|DF(z_j)u_j\|},
    \qquad
    j=-N,\dots,-1.
    \]
    Then \(u_0^{(N)}\) is used as a finite-time approximation of \(E_{z_0}^u\).

    \item Let \(z_N=F^N(z_0).\) 
    Choose any nonzero vector \(s_N\in T_{z_N}\mathbb R^2\) which is not in
    \(E_{z_N}^u\), and define
    \[
    s_{j-1}
    =
    \frac{DF(z_{j-1})^{-1}s_j}{\|DF(z_{j-1})^{-1}s_j\|},
    \qquad
    j=N,\dots,1.
    \]
    Then \(s_0^{(N)}\) is used as a finite-time approximation of \(E_{z_0}^s\).
    
\end{enumerate}
 
\end{definition}

\begin{remark}[The two right singular vectors are not used for the angle]
The singular values of $A_n(z)$ are used for finite-time Lyapunov exponents. However, the two singular vectors of \(A_n(z)\) are orthogonal because they are eigenvectors of the symmetric matrix \(A_n(z)^TA_n(z),\) which cannot measure
the true angle between \(E_z^s\) and \(E_z^u\). 
\end{remark}

\subsubsection{Angle diagnostic}

Once finite-time approximations $u_0^{(N)}$ and $s_0^{(N)}$
have been obtained, the finite-time stable--unstable angle at \(z_0\) is defined by
\[
\theta^{(N)}(z_0)
=
\arccos
\left(
\left|
\left\langle
\frac{s_0^{(N)}}{\|s_0^{(N)}\|},
\frac{u_0^{(N)}}{\|u_0^{(N)}\|}
\right\rangle
\right|
\right)
\in [0,\pi/2].
\]The absolute value is used because the directions are lines rather than oriented vectors.

For a uniformly hyperbolic invariant set, the stable and unstable directions are
uniformly transverse, so one expects
\[
\inf_{z\in\Lambda}\theta(z)>0.
\]
Therefore, very small values of
\(
\theta^{(N)}(z)
\)
along a numerical orbit suggest near-tangencies or loss of uniformity in the
finite-time hyperbolic structure.







\subsection{Numerical algorithm and results}
\label{subsec:numerical-algorithm-results}



\subsubsection{Algorithmic summary}
\label{subsubsec:algorithmic-summary}

The numerical pipeline consists of three steps: generating a long orbit, computing
finite-time Lyapunov exponents, and computing finite-time stable--unstable angles.

\begin{algorithm}[H]
\caption{Orbit generation}
\KwIn{Initial point \(z_0\), burn-in time \(N_{\mathrm{burn}}\), total length \(N_{\mathrm{tot}}\)}
\KwOut{Orbit \(z_0,\ldots,z_{N_{\mathrm{tot}}}\) after burn-in}

\For{\(j=0,\ldots,N_{\mathrm{burn}}+N_{\mathrm{tot}}-1\)}{
    \(z_{j+1}\leftarrow F(z_j)\)
}
Discard \(z_0,\ldots,z_{N_{\mathrm{burn}}-1}\)\;
Relabel the remaining orbit as \(z_0,\ldots,z_{N_{\mathrm{tot}}}\)\;
\Return{\(z_0,\ldots,z_{N_{\mathrm{tot}}}\)}
\end{algorithm}

\begin{algorithm}[H]
\caption{Finite-time Lyapunov exponents}
\KwIn{Orbit \((z_i)\), window length \(n\)}
\KwOut{\(\lambda_1^{(n)}(z_i),\lambda_2^{(n)}(z_i)\)}

\For{each sampled point \(z_i\)}{
    \(A\leftarrow I_2\)\;
    \For{\(j=0,\ldots,n-1\)}{
        \(A\leftarrow DF(z_{i+j})A\)
    }
    \((\sigma_1,\sigma_2)\leftarrow \operatorname{singular\_values}(A)\), with \(\sigma_1\geq\sigma_2\)\;
    \(\lambda_1^{(n)}(z_i)\leftarrow n^{-1}\log\sigma_1\)\;
    \(\lambda_2^{(n)}(z_i)\leftarrow n^{-1}\log\sigma_2\)\;
}
\Return{\(\lambda_1^{(n)}(z_i),\lambda_2^{(n)}(z_i)\)}
\end{algorithm}

\begin{algorithm}[H]
\caption{Finite-time stable--unstable angle}
\KwIn{Orbit \((z_i)\), propagation length \(N\)}
\KwOut{\(\theta^{(N)}(z_i)\)}

\For{each sampled point \(z_i\)}{
    Choose nonzero vectors \(u_{i-N}\) and \(s_{i+N}\)\;

    \For{\(j=i-N,\ldots,i-1\)}{
        \(u_{j+1}\leftarrow DF(z_j)u_j\)\;
        \(u_{j+1}\leftarrow u_{j+1}/\|u_{j+1}\|\)\;
    }

    \For{\(j=i+N,\ldots,i+1\)}{
        \(s_{j-1}\leftarrow DF(z_{j-1})^{-1}s_j\)\;
        \(s_{j-1}\leftarrow s_{j-1}/\|s_{j-1}\|\)\;
    }

    \(\theta^{(N)}(z_i)\leftarrow
    \arccos\!\left(|\langle u_i,s_i\rangle|\right)\)\;
}
\Return{\(\theta^{(N)}(z_i)\)}
\end{algorithm}


\subsubsection{Results}
\label{subsubsec:numerical-results}

We now apply the numerical diagnostics to the classical Hénon map. The orbit was generated from \(z_0=(0,0)\), with \(10^4\) burn-in iterates
discarded and \(10^5\) points retained. The finite-time Lyapunov exponents were
computed at \(5000\) sampled base points on the orbit.

\begin{enumerate}
    \item \textbf{Finite-time Lyapunov exponents.}
    
    The largest finite-time Lyapunov exponent was computed from the largest
    singular value of \(D(F^n)(z)\). The second exponent was obtained from
    \(
    \lambda_1^{(n)}(z)+\lambda_2^{(n)}(z)
    =
    \log |B|
    =
    \log 0.3.\)
    A long-orbit repeated-normalisation estimate gives \(\lambda_1\approx 0.4214\) and \(\lambda_2\approx -1.6254
    \).
    
    \begin{figure}[H]
    \centering

    \begin{subfigure}[t]{0.49\textwidth}
        \centering
        \begin{minipage}[c][0.27\textheight][c]{\textwidth}
            \centering
            \includegraphics[width=\textwidth,height=0.27\textheight,keepaspectratio]
            {numerical graphs/hist_lambda1_compare_n.png}
        \end{minipage}
        \caption{\centering Distributions of \(\lambda_1^{(n)}(z)\) for
        \(n=50,100,200,500\).}
        \label{fig:ftle-distributions}
    \end{subfigure}
    \hfill
    \begin{subfigure}[t]{0.49\textwidth}
        \centering
        \begin{minipage}[c][0.27\textheight][c]{\textwidth}
            \centering
            \includegraphics[width=\textwidth,height=0.27\textheight,keepaspectratio]
            {numerical graphs/mean_std_lambda1_vs_n.png}
        \end{minipage}
        \caption{\centering Mean and standard deviation of \(\lambda_1^{(n)}(z)\).}
        \label{fig:ftle-mean-std}
    \end{subfigure}

    \caption{Finite-time Lyapunov exponent diagnostics.}
    \label{fig:ftle-distribution-convergence}
    \end{figure}

    

    
    Figure~\ref{fig:ftle-distribution-convergence} shows how
    \(\lambda_1^{(n)}(z)\) depends on both the base point \(z\) and the window length
    \(n\). In (a), shorter windows give broader distributions, while larger
    \(n\) concentrates the values near the global estimate
    \(\lambda_1\approx 0.4214\). (b) shows the same convergence quantitatively:
    from \(n=50\) to \(n=500\), the mean decreases from \(0.4296\) to \(0.4223\),
    and the standard deviation decreases from \(0.0518\) to \(0.0163\). Thus
    finite-time expansion is locally nonuniform, while longer windows average these
    fluctuations toward the asymptotic Lyapunov exponent. This motivates the angle
    diagnostic below, which tests the geometric separation of stable and unstable
    directions.
    


    \item \textbf{Finite-time stable--unstable angles.}
    
    I next tested the geometric part of hyperbolicity. For a uniformly hyperbolic
    set, the stable and unstable directions are uniformly transverse, so their angle
    is bounded away from zero. For the classical Hénon parameter, however, we do not
    assume that a uniformly hyperbolic splitting exists. Rigorous non-uniform
    hyperbolicity results are known for certain Hénon parameter sets, such as the
    Benedicks--Carleson parameters, but the classical value \(B=-0.3\) is outside
    the small-dissipation regime where these results directly apply
    \cite{BenedicksCarleson1991,BenedicksYoung1993,GaliasTucker2015}. Thus the
    directions computed here are interpreted as finite-time diagnostics, not as a
    proof of a globally defined splitting.
    
    For each sampled point \(z_i\), we approximate the unstable and stable directions and the finite-time angle as in 4.1.
    
    Figure~\ref{fig:angle-summary-tail} shows the angle diagnostic. In (a),
    the angle statistics stabilise rapidly as \(N\) increases. From about \(N=10\)
    onwards, the \(1\%\) quantile, \(5\%\) quantile and median are essentially
    unchanged. This suggests that the computed directions are not dominated by the
    arbitrary choice of initial tangent vectors. We therefore use \(N=100\) as a
    representative propagation length.
    
    For \(N=100\), the median angle is approximately \(53.8^\circ\), so most sampled
    points have well-separated finite-time directions. However, (b) shows a
    small but visible lower tail near zero. In our sample, approximately \(0.44\%\)
    of points have angle below \(1^\circ\), \(0.84\%\) have angle below \(2^\circ\),
    and \(1.96\%\) have angle below \(5^\circ\). These small finite-time angles are
    interpreted as near-tangencies in the direction diagnostic.
    
 \begin{figure}[H]
    \centering

    \begin{subfigure}[t]{0.49\textwidth}
        \centering
        \begin{minipage}[c][0.27\textheight][c]{\textwidth}
            \centering
            \includegraphics[width=\textwidth,height=0.27\textheight,keepaspectratio]
            {numerical graphs/angle_summary_vs_N.png}
        \end{minipage}
        \caption{\centering Angle quantiles as the propagation length \(N\) increases.}
        \label{fig:angle-summary}
    \end{subfigure}
    \hfill
    \begin{subfigure}[t]{0.49\textwidth}
        \centering
        \begin{minipage}[c][0.27\textheight][c]{\textwidth}
            \centering
            \includegraphics[width=\textwidth,height=0.27\textheight,keepaspectratio]
            {numerical graphs/angle_lower_tail_cdf.png}
        \end{minipage}
        \caption{\centering Lower tail of \(\theta^{(100)}(z)\).}
        \label{fig:angle-lower-tail}
    \end{subfigure}

    \caption{Finite-time stable--unstable angle diagnostics. (a) checks
    convergence of the direction approximation; (b) focuses on the
    small-angle lower tail relevant to uniform transversality.}
    \label{fig:angle-summary-tail}
    \end{figure}

    Figure~\ref{fig:small-angle-spatial} separates near-tangency candidates
    \(\theta^{(100)}(z)<1^\circ\) from the broader lower tail
    \(1^\circ\leq\theta^{(100)}(z)<5^\circ\). The marked points are not spread
    uniformly over the sampled attractor; they appear mainly near bending or
    branch-close regions, especially near the right-hand turning part. This gives a
    qualitative indication that the small-angle events may be linked to the folding
    geometry of the attractor.
    
    \begin{figure}[H]
        \centering
        \includegraphics[width=0.62\textwidth]{numerical graphs/small_and_very_small_angle_points_N100.png}
        \caption{\centering Small-angle points on the sampled Hénon attractor for \(N=100\).}
        \label{fig:small-angle-spatial}
    \end{figure}
    
    Overall, the angle diagnostic does not provide evidence for a uniform lower
    bound on the finite-time stable--unstable angle. Most sampled points have
    well-separated directions, but a small lower tail of very small angles remains.
    This suggests possible non-uniform finite-time behaviour, but does not prove
    non-uniform hyperbolicity for the classical Hénon parameter.
    
    \item \textbf{Comparison between stretching and angle diagnostics.}
    
    The two diagnostics test different aspects of hyperbolic behaviour. The
    finite-time Lyapunov exponents measure stretching rates, while the angle
    diagnostic measures the separation between finite-time expanding and contracting
    directions.
    
    The Lyapunov computation gives
    \(
    \lambda_1>0, \lambda_2<0,
    \)
    showing average expansion and contraction along the sampled orbit. However,
    \(\lambda_1^{(n)}(z)\) varies across orbit segments, and the angle diagnostic
    also shows a small lower tail of near-tangency candidates. Thus the numerical
    results suggest some possible hyperbolic behaviour, but they can't indicate the uniformity of the hyperbolicity.
    

    
\end{enumerate}



\section{Conclusion}
\label{sec:conclusion}





\begin{thebibliography}{99}

\bibitem{Henon1976}
M. Hénon,
\emph{A two-dimensional mapping with a strange attractor},
Communications in Mathematical Physics, 50 (1976), 69--77.

\bibitem{BenedicksCarleson1991}
M. Benedicks and L. Carleson,
\emph{The dynamics of the Hénon map},
Annals of Mathematics, 133(1) (1991), 73--169.

\bibitem{GaliasTucker2015}
Z. Galias and W. Tucker,
\emph{Is the Hénon attractor chaotic?},
Chaos, 25(3) (2015), 033102.

\bibitem{Robinson1999}
C. Robinson,
\emph{Dynamical Systems: Stability, Symbolic Dynamics, and Chaos},
2nd ed., CRC Press, Boca Raton, 1999.

\bibitem{Barreira2012}
L. Barreira,
\emph{Ergodic Theory, Hyperbolic Dynamics and Dimension Theory},
Universitext, Springer, Berlin, 2012.

\bibitem{ColoniusKliemann2014}
F. Colonius and W. Kliemann,
\emph{Dynamical Systems and Linear Algebra},
Graduate Studies in Mathematics, Vol. 158, American Mathematical Society,
Providence, RI, 2014.

\bibitem{Sandri1996}
M. Sandri,
\emph{Numerical calculation of Lyapunov exponents},
The Mathematica Journal, 6(3) (1996), 78--84.

\end{thebibliography}



\end{document}
